\documentclass[reqno,10pt,letterpaper]{amsart}
\usepackage{mathtools,amssymb,microtype}
\usepackage[letterpaper,textwidth=16.25cm,hcentering,top=0.76in,bottom=0.82in,headheight=10pt,headsep=14pt,footskip=22pt]{geometry}
\usepackage{xcolor}
\definecolor{citeblue}{RGB}{0,70,140}
\usepackage[colorlinks=true,citecolor=citeblue,linkcolor=citeblue,urlcolor=black,pdfauthor={Wenrui Huang},pdftitle={Recent and Planned Research}]{hyperref}
\numberwithin{equation}{section}
\setlength{\emergencystretch}{1em}
\title{Recent and Planned Research}
\author{Wenrui Huang}
\date{}
\begin{document}
\maketitle

I study the long-time behavior of solutions to kinetic and dispersive equations. In a collisionless plasma, particles can approach limiting velocities while drifting logarithmically, or by a power of time, from their predicted positions. The displacement depends on the field generated by the particles themselves. I have studied this problem near vacuum in the whole space and in domains with a conducting wall. Near a spatially inhomogeneous equilibrium, some particles remain trapped, and the question becomes which part of a perturbation can damp and which part changes the equilibrium. 

\section{Scattering for particles and waves}

\subsection{Long-range interactions.}
Consider a three-dimensional particle distribution with interaction potential $\lambda |x|^{-\alpha}$. Along dispersing trajectories, the force is of order $t^{-1-\alpha}$, giving a velocity correction of order $t^{-\alpha}$. Its time integral is bounded for $\alpha>1$, logarithmic at the Coulomb exponent $\alpha=1$, and of order $t^{1-\alpha}$ for $0<\alpha<1$. With Hyunwoo Kwon, I proved small-data scattering and constructed wave operators for the Vlasov--Riesz equation in regimes on both sides of the Coulomb exponent~\cite{HK24}. On the long-range side, a kinetic pseudoconformal transformation sends infinite time to a finite endpoint. A subsequent symplectic change of variables removes the leading singularity there. We can then both extract a limiting profile from initial data and construct the solution that realizes a prescribed profile.

For interactions that decay more slowly, correcting the trajectory changes the density at the next order, which in turn changes the force. In joint work with Mengyi Xie, I studied the same difficulty for the Hartree equation and for a coupled Vlasov--Hartree system~\cite{HX26a,HX26b}. In the coupled system the particle density generates the potential acting on the wave, and the wave density generates the force on the particles. For small, regular, localized data in three dimensions and $1/2<\alpha<2$, we identified the scattering regimes of both components. At $\alpha=1$ the limiting wave density determines a logarithmic correction to the particle trajectories. The limiting particle density determines a logarithmic phase for the wave. The proof combines Hamiltonian characteristics for the particles with Galilean vector fields for the wave.

In the strongly long-range regime $0<\alpha\leq 1/2$, we found a common action--density structure for the kinetic equation, the dispersive equation, and their coupling~\cite{HX26b}. Along rays $x=ty$, a nonlinear action $\theta_k(t,y)$ modulates a wave profile by a phase or translates a kinetic profile in momentum by $\nabla_y\theta_k$. Writing $m_k$ for the rescaled spatial density of component $k$, the leading system takes the form
\begin{equation}\label{eq:action}
 \partial_t\theta_k-\frac{|\nabla_y\theta_k|^2}{2t^2}=t^{-\alpha}\varphi_j,\qquad
 \partial_t m_k-\frac{1}{t^2}\nabla_y\!\cdot(m_k\nabla_y\theta_k)
 =-\frac{1}{t^2}\nabla_y\!\cdot\mathcal J_k .
\end{equation}
Here $\varphi_j$ is the regularized potential generated by the source density and $\mathcal J_k$ is a residual current. In a single equation $j=k$, while in the coupled system the other component supplies the source. The quadratic term in \eqref{eq:action} is of order $t^{-2\alpha}$ and has a divergent time integral throughout this range. Iterating the action and density equations together yields every nonintegrable correction from the limiting densities. The $n$th correction grows like $t^{1-n\alpha}$ for $n\alpha<1$ and like $\log t$ at $n\alpha=1$. This gives small-data global existence and forward modified scattering in all three models, even as the number of necessary corrections grows without bound as $\alpha\downarrow0$.

For example, at $\alpha=1/2$ the first momentum shift is of order $t^{1/2}$ and the second is logarithmic. Both are needed for convergence. More generally, write $L_\beta(t)=\int_1^t \tau^{-\beta}\,d\tau$ and $N=\lfloor1/\alpha\rfloor$. For the kinetic equation the argument of the limiting profile contains the corrected position
\[
 x-tv-\sum_{n=1}^{N} L_{n\alpha}(t)\,\nabla\theta_n(x/t).
\]
For the wave, the recursion gives phase coefficients. In the coupled problem the two expansions are determined jointly.

\subsection{External fields and the inverse scattering problem.}
An external force can change the reference flow itself. For Vlasov--Poisson with the repulsive harmonic potential $V(x)=-|x|^2/2$, Kwon and I introduced a hyperbolic lens transform and constructed wave operators~\cite{HK26}. The transform sends $t$ to $s=\tanh t$, preserves the Hamiltonian formulation, and gives the self-consistent force a coefficient $(1-s^2)^{(d-4)/2}$. Its integrability at $s=1$ explains why small solutions scatter along the linear harmonic flow for $d\geq3$, while $d=2$ requires a logarithmic correction in the lens variables. In dimension one the coefficient is more singular. In my single-author work~\cite{H26}, I removed it by a momentum correction satisfying a forced Burgers equation. For repulsive self-interactions this correction is monotone. The resulting bounds give modified scattering for initial data of arbitrary size with one derivative and velocity weights. I also constructed modified wave operators for both signs of the interaction. In a separate dispersive problem, I constructed modified wave operators for the Hartree equation with a repulsive external Coulomb potential and small, localized scattering data~\cite{H24}.

\section{The Vlasov--Poisson system with a conducting wall}

With Benoit Pausader and Masahiro Suzuki, I studied the three-dimensional repulsive Vlasov--Poisson system in an unbounded convex domain $D$~\cite{HPS24}. Perfect conductivity imposes the Dirichlet condition $\phi|_{\partial D}=0$ on the electrostatic potential, while absorption removes particles that reach the wall. The geometry affects both which particles survive and the field they experience after traveling far from the boundary.

For small, regular, localized data satisfying the boundary compatibility conditions, we obtained global classical solutions in a class of $C^3$ convex domains. Placing the origin on the boundary, define $D_\infty=\operatorname{Int}\bigcap_{s>0}s^{-1}D$. The mass carried by velocities outside $\overline{D_\infty}$ tends to zero. If $D_\infty$ is empty, the total surviving mass tends to zero. Otherwise, particles with limiting velocities in $D_\infty$ exhibit logarithmically modified scattering. The limiting density and the Dirichlet Green function of $D_\infty$ together determine the leading electric field. In particular, the coefficient of the trajectory correction still records the geometry of the wall.

Near the boundary, we extend the potential across the wall and propagate a quantity measuring proximity to tangential impact. At the scale of an escaping particle, we decompose the Green function by spatial scale and pass from $t^{-1}D$ to the asymptotic cone. The potential and field decay as $t^{-1}$ and $t^{-2}$, respectively. Along an escaping ray the leading field is $t^{-2}E_\infty(v)$. Its tail produces a $t^{-1}$ velocity correction, whose time integral gives the logarithmic displacement.

\section{Linear stability and instability of small BGK waves}

A Bernstein--Greene--Kruskal (BGK) wave is an inhomogeneous stationary solution of the one-dimensional periodic Vlasov--Poisson system. For an energy distribution $F$, the particle Hamiltonian is $H=v^2/2-\phi(x)$. With $\min\phi=0$, the energy level $H=0$ separates particles trapped in a potential well from those that pass through it. The oscillation period changes sharply near this separatrix. In addition, varying $F$ produces a family of equilibria of infinite dimension, even after fixing the period and position of the well.

With Dongfen Bian, Emmanuel Grenier, and Benoit Pausader, I studied small one-well BGK waves bifurcating near marginal homogeneous equilibria~\cite{BGHP26}.  We obtain a criterion in terms of the dielectric function of the limiting homogeneous profile, or germ, $\mu_0(v)=F_0(v^2/2)$. Under the generic one-well hypotheses, including $\mu_0''(0)\ne0$ and the absence of nonzero neutral frequencies, a zero of the dielectric function in the open lower half-plane gives an exponentially growing mode. If there is no such zero, the wave is linearly stable up to infinite-dimensional modulation.

We use action--angle coordinates and prove quantitative resolvent estimates for the dielectric operator. Spatial Fourier modes are coupled by the inhomogeneous equilibrium, and the inverse operator must be understood near the frequencies of the trapped particles. For wave amplitude $\varepsilon$, the leading potential has the shape of a nonlinear pendulum. For $|\vartheta|\lesssim\sqrt{\varepsilon}$ the dispersion operator has the singular expansion
\[
 \mathbf M(\vartheta)=\mathbf M_h(\vartheta)+
 F'(0^+)\sqrt{\varepsilon}\,\mathbf M^{\rm lead}_{\varepsilon}
     (\vartheta/\sqrt{\varepsilon})+O_{L^2\to L^2}(\varepsilon).
\]
Here $\mathbf M_h$ is diagonal in the spatial Fourier modes of the associated homogeneous profile. The leading correction depends on the orbits of the given wave, and comparison with the pendulum gives uniform estimates on the trapping scale $|\vartheta|\sim\sqrt{\varepsilon}$. A kernel in the open lower half-plane detects a growing mode. Under the hypotheses above its growth rate tends, as the wave amplitude and distance to the germ tend to zero, to the maximal growth rate of the germ.

Translation gives an exact zero mode, while variations of the well amplitude enter at order $\sqrt{\varepsilon}$. Isolating these directions allows us to estimate the dynamical response.

On the stable side, admissible zero-charge perturbations admit the decomposition
\begin{equation}\label{eq:mod}
 f=f^{\rm bgk}+f^{\rm gal}+f^{\rm odd}+f^{\rm dyn},\qquad
 \int_0^\infty \|E^{\rm dyn}(t)\|_{L_x^2}^2\,dt<\infty .
\end{equation}
The first term varies the energy distribution within the stationary manifold. The second accounts for translation and Galilean drift. Its potential has the form $(a+bt)\partial_x\phi$, so a linearly stable wave can acquire a linearly drifting center. The third describes an imbalance between right- and left-moving passing particles that can carry current while generating no electric field.

\section{Planned research}

\subsection{Quantitative damping and nonlinear BGK dynamics.}
My immediate objective is to turn \eqref{eq:mod} into a pointwise decay estimate. After projecting out the stationary and symmetry modes, the field can be recovered by inverting the frequency-dependent dielectric operator. I plan to estimate its frequency derivatives near the trapping scale and examine the contribution from particles whose orbits approach the separatrix. This should identify how the decay rate depends on regularity of the perturbation and on the wave amplitude. One question is whether stronger regularity assumptions can yield exponential damping of the dynamical field.

The nonlinear problem also requires a choice of reference equilibrium over time. Translation and velocity alone cannot parametrize the stationary family. I plan to track the energy distribution as part of the modulation and use quantitative field decay to control the feedback of the dynamical perturbation on that distribution. In the unstable regime, my goal is to control the nonlinear error over the time required for a growing linear mode to reach a fixed size. A separate problem is the degenerate threshold $\mu_0''(0)=0$. Here the leading trapping contribution used in the generic argument vanishes, and I want to identify the first surviving term in the dispersion operator.

\subsection{Specular reflection and grazing trajectories.}
I plan to keep the conducting boundary condition for the potential and replace absorption by specular reflection for the particles. With unit normal $n(x)$, the outgoing velocity is $v-2(v\cdot n)n$. The absorbing problem ends a characteristic at its first impact. With reflection, a near-tangential trajectory can return many times, and the field generated between impacts changes the angle of the next collision. My first objective is to determine whether our grazing estimates can be propagated across successive reflections with sufficient control of derivatives of the characteristic flow. I would then study global regularity and the contribution of repeatedly reflected particles to the asymptotic field.

\subsection{External potentials and confinement.}
I would like to identify radial external potentials for which one can construct asymptotic coordinates for the self-consistent dynamics. For spherically symmetric solutions, I plan to start from the integrable radial motion of the external Hamiltonian at fixed angular momentum and determine the correction due to the evolving density. For nonradial perturbations the self-consistent force can also change angular momentum, and this effect has to be included in the coordinates. The attractive harmonic potential gives a complementary test in which the unperturbed trajectories are periodic. There I want to use the lens transform to study how the self-consistent force accumulates over successive periods and how that accumulation affects phase mixing. In the harmonic flow all unperturbed orbits have the same period. For a general radial potential the frequency can vary with the orbit, allowing nearby trajectories to separate in phase. I aim to quantify how that variation affects field decay.

\begingroup
\small
\renewcommand{\refname}{References}
\begin{thebibliography}{9}
\setlength{\itemsep}{1pt}\setlength{\parskip}{0pt}\setlength{\parsep}{0pt}
\bibitem{HK24} W. Huang and H. Kwon. \href{https://doi.org/10.1007/s00023-026-01735-7}{Scattering of the Vlasov--Riesz System in the Three Dimensions}. \emph{Ann. Henri Poincar\'e} (2026).
\bibitem{HX26a} W. Huang and M. Xie. Long-time dynamics of Vlasov--Hartree systems across the Coulomb threshold. \href{https://arxiv.org/abs/2608.04402}{arXiv:2608.04402}.
\bibitem{HX26b} W. Huang and M. Xie. A common structure for modified scattering: Vlasov--Riesz, Hartree, and their coupling. \href{https://arxiv.org/abs/2610.02643}{arXiv:2610.02643}.
\bibitem{HK26} W. Huang and H. Kwon. Scattering map for the Vlasov--Poisson system with a repulsive harmonic potential. \href{https://arxiv.org/abs/2602.21344}{arXiv:2602.21344}.
\bibitem{H26} W. Huang. Asymptotic behavior of the one-dimensional Vlasov--Poisson system with a repulsive harmonic potential. \href{https://arxiv.org/abs/2610.02551}{arXiv:2610.02551}.
\bibitem{H24} W. Huang. Modified wave operators for the Hartree equation with repulsive Coulomb potential. \href{https://arxiv.org/abs/2411.02737}{arXiv:2411.02737}.
\bibitem{HPS24} W. Huang, B. Pausader, and M. Suzuki. \href{https://doi.org/10.1007/s00220-026-05731-0}{The Vlasov--Poisson System with a Perfectly Conducting Wall: Convex Domains}. \emph{Comm. Math. Phys.} \textbf{407}, 219 (2026).
\bibitem{BGHP26} D. Bian, E. Grenier, W. Huang, and B. Pausader. Stability and instability of small BGK waves. \href{https://arxiv.org/abs/2601.10030}{arXiv:2601.10030}.
\end{thebibliography}
\endgroup
\end{document}
